%======================================================================
%  style.tex -- page layout and typography
%  Visual language taken from the reciprocity-gap notes (Chapter 3):
%  sans-serif bold headings, light-gray boxes, exercises with
%  solutions in small type.
%======================================================================
\usepackage[margin=2.4cm,headheight=14pt]{geometry}
% Latin Modern with T1 when installed (standard TeX Live); plain Computer
% Modern otherwise, so that the file compiles on minimal installations.
\IfFileExists{lmodern.sty}{\usepackage{lmodern}\usepackage[T1]{fontenc}}{}
\usepackage[utf8]{inputenc}
\usepackage{amsmath,amssymb,amsthm,bm,mathtools}
\usepackage{xcolor,graphicx,enumitem,booktabs,array,multirow}
\usepackage{tikz}
\usetikzlibrary{arrows.meta,calc,positioning,shapes.geometric,
  decorations.pathreplacing,decorations.markings,decorations.pathmorphing,
  patterns}
\usepackage{pgfplots}\pgfplotsset{compat=1.18}  % 3D sketches of Chapter 2 (figures/tikz)
\usepackage[most]{tcolorbox}
\usepackage[font=small,labelfont={sf,bf}]{caption}
\usepackage{listings}
\usepackage{algorithm}      % CMAME-style pseudocode in the exercise parts
\usepackage{algpseudocode}
\usepackage{microtype}
\usepackage{etoolbox}
\usepackage{fancyhdr}
\usepackage{imakeidx}
\makeindex[intoc,columns=2,options=-s preamble/index.ist]
\usepackage[explicit]{titlesec}
\usepackage{hyperref}
\usepackage{xurl}           % URLs break at any character (Ch. 7 links)
\hypersetup{colorlinks,linkcolor=blue!50!black,citecolor=blue!50!black,
  urlcolor=blue!50!black,pdftitle={Nonlinear Problems in Mechanics -- Lecture notes},
  pdfauthor={Andrei Constantinescu}}

%---- headings: sans-serif bold, as in the RG notes ---------------------
\titleformat{\chapter}[display]
  {\sffamily\bfseries}
  {\large\color{gray!70!black}\MakeUppercase{\chaptertitlename}~\thechapter}
  {4pt}
  {\titlerule[0.6pt]\vspace{6pt}\Huge #1}
  [\vspace{2pt}]
\titleformat{name=\chapter,numberless}[display]
  {\sffamily\bfseries}{}{0pt}{\titlerule[0.6pt]\vspace{6pt}\Huge #1}
\titlespacing*{\chapter}{0pt}{-10pt}{28pt}
\titleformat{\section}{\sffamily\bfseries\Large}{\thesection}{0.8em}{#1}
\titleformat{name=\section,numberless}{\sffamily\bfseries\Large}{}{0pt}{#1}
\titleformat{\subsection}{\sffamily\bfseries\large}{\thesubsection}{0.8em}{#1}
\titleformat{name=\subsection,numberless}{\sffamily\bfseries\large}{}{0pt}{#1}
\titleformat{\subsubsection}{\sffamily\bfseries}{\thesubsubsection}{0.8em}{#1}
\titleformat{name=\subsubsection,numberless}{\sffamily\bfseries}{}{0pt}{#1}
\titleformat{\paragraph}[runin]{\sffamily\bfseries}{}{0pt}{#1}
\titlespacing*{\paragraph}{0pt}{1.6ex plus .5ex minus .2ex}{0.6em}

%---- running heads -------------------------------------------------------
\pagestyle{fancy}
\fancyhf{}
\renewcommand{\headrulewidth}{0.4pt}
\fancyhead[LE]{\sffamily\small\leftmark}
\fancyhead[RO]{\sffamily\small\rightmark}
\fancyfoot[C]{\sffamily\small\thepage}
\renewcommand{\chaptermark}[1]{\markboth{\thechapter.\ #1}{}}
\renewcommand{\sectionmark}[1]{\markright{\thesection\ \ #1}}
\fancypagestyle{plain}{\fancyhf{}\renewcommand{\headrulewidth}{0pt}%
  \fancyfoot[C]{\sffamily\small\thepage}}

%---- boxes: light gray background, bold sans-serif title -----------------
% A \label at the top of a gbox points to the box: \ref gives the number of the
% current section, \nameref the title of the box. Key statements of the
% plasticity chapters are cited this way ("the projection property,
% Section~\ref{...}") instead of Theorem/Proposition numbers.
\makeatletter
\newcommand{\gboxname}[1]{\phantomsection\def\@currentlabelname{#1}}
\makeatother
\newtcolorbox{gbox}[1]{enhanced,breakable,colback=gray!12,colframe=gray!12,
  boxrule=0pt,arc=2pt,left=8pt,right=8pt,top=6pt,bottom=6pt,
  title={#1},fonttitle=\sffamily\bfseries,coltitle=black,colbacktitle=gray!12,
  toptitle=3pt,bottomtitle=0pt,titlerule=0pt,before upper={\gboxname{#1}}}
% untitled variant, used for take-home messages and short summaries
\newtcolorbox{gbox*}{enhanced,breakable,colback=gray!12,colframe=gray!12,
  boxrule=0pt,arc=2pt,left=8pt,right=8pt,top=6pt,bottom=6pt}

%---- theorem-like statements in the same sans-serif voice ---------------
\newtheoremstyle{sfthm}{\topsep}{\topsep}{\itshape}{}{\sffamily\bfseries}{.}{0.5em}{}
\theoremstyle{sfthm}
\newtheorem{theorem}{Theorem}[chapter]
\newtheorem{proposition}[theorem]{Proposition}
\newtheorem{lemma}[theorem]{Lemma}
\newtheorem{corollary}[theorem]{Corollary}
\newtheoremstyle{sfdef}{\topsep}{\topsep}{\normalfont}{}{\sffamily\bfseries}{.}{0.5em}{}
\theoremstyle{sfdef}
\newtheorem{definition}[theorem]{Definition}
\newtheorem{remark}[theorem]{Remark}
\newtheorem{example}[theorem]{Example}
\renewcommand{\proofname}{\sffamily\bfseries Proof}

%---- exercises: sans-serif heading, solution in small type ---------------
\newcounter{exo}[chapter]
\renewcommand{\theexo}{\thechapter.\arabic{exo}}
\newcommand{\exercise}[1]{\par\medskip\noindent\refstepcounter{exo}%
  {\sffamily\bfseries Exercise~\theexo\quad#1.}\ }
\newcommand{\solsize}{\scriptsize}   % change to \tiny for smaller solutions
\newenvironment{solution}{\par\smallskip\begingroup\solsize\noindent
  {\sffamily\bfseries Solution.}\ }{\par\endgroup\medskip}

%---- code listings (Python), quiet gray frame in the box colour ----------
\lstdefinestyle{pystyle}{language=Python,
  basicstyle=\ttfamily\scriptsize,
  keywordstyle=\color{blue!60!black},commentstyle=\color{gray},
  stringstyle=\color{green!45!black},
  backgroundcolor=\color{gray!6},frame=none,
  breaklines=true,showstringspaces=false,tabsize=4,
  xleftmargin=0.6em,aboveskip=4pt,belowskip=4pt}
\lstset{style=pystyle}

%---- code listings (gibiane, the language of Cast3M) ---------------------
% Comments are the lines with a * in the first column (fixed-column comment
% of listings); keywords are case-insensitive, as in gibiane. Use with
% \lstinputlisting[style=gibstyle]{file.dgibi}.
\lstdefinelanguage{gibiane}{%
  sensitive=false,
  morecomment=[f][\color{gray}][0]{*},
  morestring=[b]',
  morekeywords={opti,debproc,finproc,repeter,fin,si,sinon,finsi,quitter,
    mess,table,prog,mots,lect,et,ega,dime,extr,exco,masq,vmis,sigm,bsig,
    reso,rigi,mode,mate,cara,bloq,depi,reac,pres,forc,char,evol,pasapas,
    droi,dall,cerc,coor,chan,zero,manu,dess,trac,abs,maxi,exp,log,flot,
    exis,poin,ipol,rela,elim,resu,nbno,sin,cos,inve,index,vale,chai,mot},
  keywordstyle=\color{blue!60!black}}
\lstdefinestyle{gibstyle}{language=gibiane,
  basicstyle=\ttfamily\scriptsize,
  stringstyle=\color{green!45!black},
  backgroundcolor=\color{gray!6},frame=none,
  breaklines=true,showstringspaces=false,columns=fixed,
  xleftmargin=0.6em,aboveskip=4pt,belowskip=4pt}

%---- TikZ styles shared by the diagrams ---------------------------------
\definecolor{cu}{RGB}{31,90,200}     % potential / displacement data S_u
\definecolor{cphi}{RGB}{20,150,60}   % flux / traction data S_phi, S_t
\tikzset{
  box/.style={rectangle,rounded corners=2pt,fill=gray!12,draw=none,
    minimum width=8.2cm,minimum height=1.0cm,align=center,inner sep=5pt,
    font=\small},
  smallbox/.style={rectangle,rounded corners=2pt,fill=gray!12,draw=none,
    minimum width=4.6cm,minimum height=1.25cm,align=center,inner sep=5pt,
    font=\small},
  tinybox/.style={rectangle,rounded corners=2pt,fill=gray!12,draw=none,
    minimum width=2.7cm,text width=2.4cm,minimum height=1.3cm,align=center,
    inner sep=4pt,font=\footnotesize},
  goal/.style={fill=green!12},
  % Tonti diagrams: kinematic quantities in blue, static quantities in green
  kin/.style={fill=cu!12,draw=cu!45,thin},
  stat/.style={fill=cphi!12,draw=cphi!45,thin},
  kinlab/.style={font=\small\sffamily,text=cu},
  statlab/.style={font=\small\sffamily,text=cphi},
  key/.style={fill=orange!14},
  arr/.style={-{Latex[length=2.3mm]},thick},
  spring/.style={decorate,decoration={coil,aspect=0,segment length=4pt,
    amplitude=3pt,pre length=0.15cm,post length=0.15cm},thick}
}

%---- framing parts of every chapter --------------------------------------
% Outline, Summary, Exercises, References: unnumbered, roman small caps,
% listed in the table of contents. \framesub is the level below (e.g. the
% parts of the exercises). A \label placed after \framepart works with
% \nameref and \pageref (not with \ref, since the part has no number).
\makeatletter
\newcommand{\framepart}[1]{%
  \par\addvspace{4.5ex plus 1ex minus .5ex}%
  \phantomsection
  \addcontentsline{toc}{section}{\protect\textsc{#1}}%
  \markright{#1}%
  \def\@currentlabelname{#1}%
  {\noindent\normalfont\Large\scshape #1\par}%
  \nobreak\vspace{0.6ex}\nobreak
  {\color{gray!55}\hrule height 0.4pt}%
  \nobreak\vspace{1.6ex plus .3ex}%
  \@afterheading}
\newcommand{\framesub}[1]{%
  \par\addvspace{3ex plus .8ex minus .3ex}%
  {\noindent\normalfont\large\scshape #1\par}%
  \nobreak\vspace{1ex plus .2ex}%
  \@afterheading}
\makeatother

%---- per-chapter bibliographies -----------------------------------------
% Each chapter ends with \begin{chapterbib}{99} ... \end{chapterbib}.
% Numbering restarts in every chapter; keys must be unique book-wide.
\makeatletter
\newenvironment{chapterbib}[1]
  {\framepart{References}%
   \small
   \list{\@biblabel{\@arabic\c@enumiv}}%
     {\settowidth\labelwidth{\@biblabel{#1}}%
      \leftmargin\labelwidth\advance\leftmargin\labelsep
      \itemsep 2pt\usecounter{enumiv}%
      \let\p@enumiv\@empty\renewcommand\theenumiv{\@arabic\c@enumiv}}%
   \sloppy\clubpenalty4000\widowpenalty4000\sfcode`\.\@m}
  {\def\@noitemerr{\@latex@warning{Empty `chapterbib' environment}}%
   \endlist}
\makeatother

%---- drafting aids (set \drafttrue/\draftfalse in main.tex) --------------
\newif\ifdraft
\newcommand{\draftnote}[1]{\ifdraft{\sffamily\footnotesize\color{red!70!black}[#1]}\fi}

\numberwithin{equation}{chapter}
\numberwithin{figure}{chapter}
\numberwithin{table}{chapter}
\numberwithin{algorithm}{chapter}
\setlist{itemsep=1pt}
